\documentclass[a4paper,11pt]{article}

% --- Paquetes ---
\usepackage[utf8]{inputenc}
\usepackage[T1]{fontenc}
\usepackage[spanish]{babel}
\usepackage{geometry}
\usepackage{booktabs}
\usepackage{enumitem}
\usepackage{xcolor}
\usepackage{hyperref}
\usepackage{fancyhdr}
\usepackage{titlesec}
\usepackage{tabularx}
\usepackage{multicol}
\usepackage{microtype}
\usepackage{tikz}

\geometry{margin=2cm, top=2.5cm, bottom=2.5cm}
\setlength{\parindent}{0pt}

% --- Colores ---
\definecolor{cmdcolor}{HTML}{2E86C1}
\definecolor{paramcolor}{HTML}{E67E22}
\definecolor{notecolor}{HTML}{7D3C98}
\definecolor{sectioncolor}{HTML}{1A5276}
\definecolor{warncolor}{HTML}{C0392B}
\definecolor{sectionbg}{HTML}{D6EAF8}

% --- Formato de secciones ---
\titleformat{\section}{\Large\bfseries\color{sectioncolor}}{}{0em}{%
  \tikz[overlay]\fill[sectionbg, rounded corners=2pt]
    (-2pt,-3pt) rectangle (\textwidth+2pt,\baselineskip+2pt);%
}
\titleformat{\subsection}{\large\bfseries\color{sectioncolor!80}}{}{0em}{}

% --- Comando para formatear comandos ---
\newcommand{\cmd}[1]{\texttt{\color{cmdcolor}\bfseries #1}}
\newcommand{\param}[1]{\texttt{\color{paramcolor}<#1>}}
\renewcommand{\bot}[0]{\texttt{[bot]}}
\newcommand{\nota}[1]{\par\smallskip{\small\color{notecolor}\textbf{Nota:} #1}\par}
\newcommand{\aviso}[1]{\par\smallskip{\small\color{warncolor}\textbf{¡Atención!} #1}\par}

% --- Encabezado y pie ---
\pagestyle{fancy}
\fancyhf{}
\fancyhead[L]{\small\color{sectioncolor}Guía rápida para DJs — Bot IRC}
\fancyhead[R]{\small\color{sectioncolor}\thepage}
\fancyfoot[C]{\small\color{gray}Febrero 2026 — El Desván Musical}
\renewcommand{\headrulewidth}{0.4pt}

% --- Documento ---
\begin{document}

% === Portada ===
\begin{center}
  {\Huge\bfseries\color{sectioncolor} Guía Rápida para DJs}\\[0.5em]
  {\Large Bot IRC — El Desván Musical}\\[1.5em]
  {\large Referencia abreviada de comandos y procedimientos}\\[0.5em]
  {\color{gray}\today}
\end{center}

\vspace{1em}

\begin{center}
\fbox{\parbox{0.9\textwidth}{
\textbf{Convenciones de este documento:}
\begin{itemize}[nosep,leftmargin=1.5em]
  \item \bot{} debe sustituirse por el nick del bot (ej.\ \texttt{c3pO}).
  \item \param{param} indica un parámetro que debes rellenar.
  \item Los comandos con \cmd{!} se envían como mensaje normal.
  \item Los comandos con \cmd{/me} se envían como acción de IRC.
  \item Los comandos marcados como \textbf{PM} se ejecutan en privado con el bot.
\end{itemize}
}}
\end{center}

\vspace{0.5em}

% ============================================================
\section{Flujo de una emisión}
% ============================================================

El flujo típico de una emisión es:

\begin{enumerate}[leftmargin=2em]
  \item \textbf{Iniciar emisión:} \cmd{!emito}
  \item \textbf{Abrir peticiones:} \cmd{/me abre peticiones} (desde el canal)
  \item \textbf{Gestionar peticiones:} ver, poner, eliminar peticiones según llegan
  \item \textbf{Cerrar peticiones:} \cmd{/me cierra peticiones}
  \item \textbf{Finalizar emisión:} \cmd{!auto}
\end{enumerate}

\nota{Si el DJ abandona el canal (PART/QUIT), el bot espera 10 minutos. Si no regresa, activa automáticamente el modo \texttt{auto}.}

% ============================================================
\section{Comandos de emisión}
% ============================================================

\begin{tabularx}{\textwidth}{lX}
\toprule
\textbf{Comando} & \textbf{Descripción} \\
\midrule
\cmd{!emito} & Registra tu nick como DJ activo. Anuncia en todos los canales ``Feliz emisión \param{nick}''. Se guarda en la configuración. \\[0.5em]
\cmd{!auto} & Finaliza la emisión. Cambia el DJ a ``Auto'', anuncia el modo automático con la URL de la radio, limpia todas las peticiones, resetea el límite a~0 y elimina los archivos de peticiones. \\
\bottomrule
\end{tabularx}

% ============================================================
\section{Peticiones}
% ============================================================

\subsection{Abrir y cerrar peticiones}

\begin{tabularx}{\textwidth}{lX}
\toprule
\textbf{Comando} & \textbf{Descripción} \\
\midrule
\cmd{/me abre peticiones} & Abre las peticiones. Solo funciona si eres el DJ activo. El mensaje debe contener alguna forma de ``abrir'' y ``petición''. El bot anuncia la apertura y la sintaxis para pedir en todos los canales. \\[0.5em]
\cmd{/me cierra peticiones} & Cierra las peticiones. El mensaje debe contener ``cerrar'' y ``petición''. La cola \textbf{no} se borra (puedes seguir usando \cmd{!suena}). \\
\bottomrule
\end{tabularx}

\nota{Los usuarios hacen peticiones con \cmd{!peticion cantante-canción-dedicatoria}. El bot te las envía por PM numeradas.}

\subsection{Gestión de la cola de peticiones}

Todos estos comandos se ejecutan \textbf{por PM} al bot.

\begin{tabularx}{\textwidth}{lX}
\toprule
\textbf{Comando} & \textbf{Descripción} \\
\midrule
\cmd{!petlista} & Muestra la lista de peticiones pendientes (numeradas). \\[0.3em]
\cmd{!suena \param{N}} & Anuncia la petición número $N$ como ``Está sonando\ldots'' en todos los canales y la elimina de la cola. \\[0.3em]
\cmd{!elimina \param{N}} & Elimina la petición número $N$ de la cola de forma silenciosa (sin anuncio público). \\[0.3em]
\cmd{!recuento} & Muestra el recuento de peticiones por usuario. \\
\bottomrule
\end{tabularx}

\subsection{Límites de peticiones por usuario}

\begin{tabularx}{\textwidth}{lX}
\toprule
\textbf{Comando} & \textbf{Descripción} \\
\midrule
\cmd{!limite \param{N}} & Establece el límite máximo de peticiones por usuario a $N$. Se anuncia en los canales. Con $N=0$ se elimina el límite. Se guarda en la configuración. \\[0.3em]
\cmd{!muestralimite} & Muestra el límite actual. \\[0.3em]
\cmd{!-\param{N} \param{nick}} & Resta $N$ al contador de peticiones de \param{nick}. Útil para anular una petición o conceder una extra. Ejemplo: \cmd{!-2 fulanito} \\
\bottomrule
\end{tabularx}

\subsection{Ignorar y readmitir usuarios}

Para gestionar usuarios problemáticos durante la emisión:

\begin{tabularx}{\textwidth}{lX}
\toprule
\textbf{Comando} & \textbf{Descripción} \\
\midrule
\cmd{!ignora \param{nick}} & Ignora las peticiones futuras de \param{nick}. \\[0.3em]
\cmd{!admite \param{nick}} & Readmite a \param{nick} (deja de ignorarlo). \\[0.3em]
\cmd{!iglista} & Muestra la lista de usuarios ignorados. \\
\bottomrule
\end{tabularx}

\aviso{La lista de ignorados es volátil: se pierde al reiniciar el bot.}

% ============================================================
\section{Control del bot}
% ============================================================

\begin{tabularx}{\textwidth}{lX}
\toprule
\textbf{Comando} & \textbf{Descripción} \\
\midrule
\cmd{/invite \param{bot} \param{canal}} & Invita al bot al canal (comando IRC est\'{a}ndar). Solo administradores. \\[0.3em]
\bot{} \cmd{sal del canal} & El bot abandona el canal actual. Es temporal: se re-une en el siguiente chequeo peri\'{o}dico o PING. Para abandonar de forma permanente, un admin debe usar \cmd{!leave}. \\[0.3em]
\bot{} \cmd{recarga} & El bot se desconecta del IRC. El crontab lo reinicia automáticamente en $\sim\!1$~minuto. Se puede añadir texto extra, ej.\ ``\texttt{c3pO salte a recargar un momento}''. \\
\bottomrule
\end{tabularx}

% ============================================================
\section{Información de la radio}
% ============================================================

\begin{tabularx}{\textwidth}{lX}
\toprule
\textbf{Comando} & \textbf{Descripción} \\
\midrule
\cmd{!radio} / \cmd{!link} & Muestra en todos los canales: enlace web de la radio, enlace de streaming, DJ actual y estado de las peticiones (abiertas/cerradas). \\
\bottomrule
\end{tabularx}

% ============================================================
\section{Entretenimiento y utilidades}
% ============================================================

Estos comandos están disponibles para todos los usuarios y pueden ser útiles durante la emisión:

\begin{tabularx}{\textwidth}{lX}
\toprule
\textbf{Comando} & \textbf{Descripción} \\
\midrule
\bot{} \texttt{refrán} & El bot cuenta un refrán. El mensaje debe contener el nick del bot y la palabra ``refrán'' (o ``refran''), en cualquier orden. Ej.: ``\textit{c3pO dinos un refrán}''. \\[0.3em]
\bot{} \texttt{chiste} & El bot cuenta un chiste. Ej.: ``\textit{Cuéntame un chiste c3pO}''. \\[0.3em]
\bot{} \texttt{horóscopo} \param{signo} & El bot da el horóscopo del día para el signo zodiacal indicado. Ej.: ``\textit{c3pO horóscopo libra}''. \\[0.3em]
\cmd{!help} / \cmd{!ayuda} & Muestra los grupos de ayuda disponibles. Con \cmd{!help \param{grupo}} se listan los comandos de ese grupo (responde por PM). Grupos visibles para DJs: \texttt{general}, \texttt{peticiones}, \texttt{dj}. \\
\bottomrule
\end{tabularx}

% ============================================================
\section{Interacción con la IA}
% ============================================================

En los canales donde la IA está activada, puedes interactuar con el bot mencionando su nick:

\begin{itemize}[leftmargin=2em]
  \item \textbf{En canal:} escribe un mensaje que contenga el nick del bot y este responderá usando inteligencia artificial. Ej.: ``\textit{c3pO ¿qué opinas de la canción?}''
  \item \textbf{Por PM:} al ser DJ autorizado, puedes enviar mensajes privados al bot. Si el mensaje no empieza por \cmd{!}, el bot lo interpreta como una consulta a la IA.
\end{itemize}

\nota{Los DJs no pueden configurar la IA (modelo, prompt, canales). Esos ajustes son exclusivos de administradores.}

% ============================================================
\section{Mensajes privados (PM)}
% ============================================================

Como DJ autorizado, el bot acepta tus mensajes privados. Los usuarios normales no pueden escribir al bot por privado.

La mayoría de los comandos de gestión (\cmd{!petlista}, \cmd{!suena}, \cmd{!elimina}, \cmd{!ignora}, etc.) se ejecutan por PM para no saturar el canal.

% ============================================================
\newpage
\section{Referencia rápida}
% ============================================================

\begin{multicols}{2}
\small

\textbf{Emisión}
\begin{itemize}[nosep,leftmargin=1em]
  \item \cmd{!emito} — iniciar emisión
  \item \cmd{!auto} — finalizar emisión
\end{itemize}

\textbf{Peticiones (canal)}
\begin{itemize}[nosep,leftmargin=1em]
  \item \cmd{/me abre peticiones}
  \item \cmd{/me cierra peticiones}
\end{itemize}

\textbf{Peticiones (PM)}
\begin{itemize}[nosep,leftmargin=1em]
  \item \cmd{!petlista} — ver cola
  \item \cmd{!suena \param{N}} — poner petición
  \item \cmd{!elimina \param{N}} — borrar petición
  \item \cmd{!recuento} — conteo por usuario
\end{itemize}

\textbf{Límites (PM)}
\begin{itemize}[nosep,leftmargin=1em]
  \item \cmd{!limite \param{N}} — fijar límite
  \item \cmd{!muestralimite} — ver límite
  \item \cmd{!-\param{N} \param{nick}} — restar
\end{itemize}

\columnbreak

\textbf{Ignorar/admitir (PM)}
\begin{itemize}[nosep,leftmargin=1em]
  \item \cmd{!ignora \param{nick}}
  \item \cmd{!admite \param{nick}}
  \item \cmd{!iglista}
\end{itemize}

\textbf{Control del bot}
\begin{itemize}[nosep,leftmargin=1em]
  \item \bot{} \texttt{sal del canal}
  \item \bot{} \texttt{recarga} — reiniciar
  \item \cmd{/invite} \param{bot} — invitar
\end{itemize}

\textbf{General}
\begin{itemize}[nosep,leftmargin=1em]
  \item \cmd{!radio} — info de la radio
  \item \bot{} \texttt{refrán}
  \item \bot{} \texttt{chiste}
  \item \bot{} \texttt{horóscopo} \param{signo}
  \item \cmd{!help} / \cmd{!ayuda}
\end{itemize}

\end{multicols}

\end{document}
